\documentclass[12pt,letterpaper]{article}
\usepackage[utf8]{inputenc}
\usepackage[spanish]{babel}
\usepackage{amsmath}
\usepackage{amsfonts}
\usepackage{amssymb}
\usepackage[margin=2.5cm]{geometry}

\begin{document}

\begin{center}
    \Large \textbf{Evaluación de Examen 2: Unidad II} \\
    \large Física para Ciencias de la Salud (005-1713) \\
    \vspace{0.5cm}
    \textbf{Estudiante:} Camila Cabrera \quad \textbf{C.I.:} 32.937.064
    \vspace{0.3cm} \\
    \large \textbf{Calificación Final: 7,5/10 puntos}
\end{center}

\vspace{0.5cm}

\section*{Análisis de Resultados}

\subsection*{1. Cinemática (Aceleración media)}
\textbf{Procedimiento y resultado:} Extrae los datos y aplica de forma correcta la fórmula de la aceleración media ($a = \frac{v_f - v_i}{t}$). La división aritmética de $0,6 / 0,15$ es correcta ($4$); sin embargo, falla en el análisis dimensional al expresar el resultado final con unidades de velocidad ($4 \text{ m/s}$) en lugar de la unidad de aceleración ($4 \text{ m/s}^2$). \\
\textbf{Veredicto:} \textbf{Parcialmente correcto} (Penalización por unidades incorrectas) (1,5/2 pts).

\subsection*{2. Dinámica (Leyes de Newton)}
\textbf{Procedimiento y resultado:} Identifica correctamente los datos. No obstante, al momento de operar comete un error conceptual grave: en lugar de aplicar el despeje analítico de la Segunda Ley de Newton para la aceleración ($a = \frac{F}{m}$), plantea erróneamente una multiplicación entre la masa y la fuerza ($25 \text{ kg} \cdot 150 \text{ N}$). Esto arroja un valor incorrecto de $3750$ acompañado de unidades equivocadas, cuando el resultado real era $6 \text{ m/s}^2$. \\
\textbf{Veredicto:} \textbf{Incorrecto} (Se otorgan puntos parciales por la identificación inicial de los datos) (0,5/2 pts).

\subsection*{3. Trabajo Mecánico}
\textbf{Procedimiento y resultado:} Muestra un procedimiento excelente. Plantea correctamente la fórmula del trabajo mecánico y sustituye los valores correspondientes ($400 \text{ N} \cdot 0,8 \text{ m}$). El cálculo matemático arroja el valor exacto de $320$, acompañado correctamente de su unidad en el Sistema Internacional (Joules). \\
\textbf{Veredicto:} \textbf{Correcto} (2/2 pts).

\subsection*{4. Energía Potencial}
\textbf{Procedimiento y resultado:} Extrae de manera estructurada los datos del enunciado y aplica a la perfección la fórmula de la energía potencial gravitatoria ($E_p = m \cdot g \cdot h$). El cálculo de $1,2 \text{ kg} \cdot 9,8 \text{ m/s}^2 \cdot 1,5 \text{ m}$ arroja el valor exacto y verificado de $17,64 \text{ J}$. \\
\textbf{Veredicto:} \textbf{Correcto} (2/2 pts).

\subsection*{5. Potencia Mecánica}
\textbf{Procedimiento y resultado:} La estudiante identifica correctamente la ecuación de la potencia mecánica ($P = \frac{W}{t}$). Sin embargo, comete un error de atención en la extracción de datos, anotando el trabajo como $72 \text{ J}$ en lugar de los $75 \text{ J}$ indicados en el enunciado. Desarrolla la operación con este dato equivocado ($72 \text{ J} / 60 \text{ s}$), obteniendo consecuentemente $1,2 \text{ W}$ en lugar del verdadero resultado de $1,25 \text{ W}$. \\
\textbf{Veredicto:} \textbf{Parcialmente correcto} (Se penaliza el error de transcripción de datos, aunque el desarrollo matemático posterior guarde coherencia con dicho error) (1,5/2 pts).

\vspace{0.5cm}
\section*{Comentarios Generales y Sugerencias}
El examen presenta un formato bastante ordenado, extrayendo metódicamente los datos en el lado izquierdo y desarrollando las operaciones a la derecha, logrando puntajes perfectos en Trabajo y Energía.
\begin{itemize}
    \item Se sugiere leer muy cuidadosamente los datos del enunciado. Un simple error de transcripción (como cambiar 75 por 72 en la pregunta 5) puede alterar todo el resultado final de un paciente o muestra clínica real.
    \item Reforzar los **despejes matemáticos**. Identificar si una variable está multiplicando o dividiendo es crucial para llegar a la respuesta física correcta, tal como ocurrió en la pregunta de las Leyes de Newton.
    \item Prestar máxima atención a las **unidades de las magnitudes**. La aceleración en el Sistema Internacional siempre lleva la unidad de tiempo al cuadrado ($\text{m/s}^2$).
\end{itemize}

\end{document}